\subsection*{Pregunta 4}
\textbf{¿Cuáles son las ventajas, para un DBA el trabajar con un SMBD, oper source?}
\addcontentsline{toc}{subsection}{Pregunta 4}

Para un administrador de bases de datos, trabajar con un SMBD open source le permite tener mayor libertad para innovar y moldear el sistema a sus propias necesidades. Aunque esto implica que el administrador debe tener un mayor conocimiento y experiencia con la infraestructura que utiliza, también tiene sus beneficios, tales como:
\begin{itemize}
    \item Reducir costos por uso de licencias. Al ser OpenSource, son de uso gratuito, lo que permite reducir el costo de implementar el proyecto y dirigirlo en beneficio del DBA (Como en capacitación).
    \item Reducir costos de mantenimiento. Los SMBD open source se sostienen a partir del conocimiento compartido entre la comunidad  (como es el caso de PostgreSQL), en las que se comparten herramientas de uso y soluciones a fallos, lo que reduce costos de mantenimiento.
    \item Sin restricciones por parte del proveedor. Al no estar limitado por lo que permite el proveedor, se tiene la capacidad de adaptar el sistema a la problemática a enfrentar. 
    \item Autonomía y transparencia en cuestiones de seguridad. Por una parte, al no existir una dependencia con el proveedor al momento de necesitar soluciones, se vuelve responsabilidad del DBA examinar y corregir cualquier vulnerabilidad en cuanto se encuentre. No obstante, esto también implica tener un mayor control sobre el uso de sus datos y la forma de utilizarlos.
\end{itemize}
\newpage
